\documentclass[11pt,a4paper]{book}
\usepackage[left=3cm,right=3cm,top=3cm,bottom=3cm]{geometry}
\usepackage[onehalfspacing]{setspace}

\usepackage{amsmath, amssymb, amsthm}

\usepackage{hyperref}
\urlstyle{same}

\usepackage{tikz,tikz-cd}

\usepackage{bm}

%\usepackage[backend=biber]{biblatex}
%\addbibresource{biblio_guillermo.bib}

\newtheorem{thm}{Theorem}[section]
\newtheorem{prop}[thm]{Proposition}
\newtheorem{conj}[thm]{Conjecture}
\newtheorem{lemma}[thm]{Lemma}
\newtheorem{corol}[thm]{Corollary}
\theoremstyle{remark} \newtheorem{rmk}[thm]{Remark}
\theoremstyle{definition} \newtheorem{defn}[thm]{Definition}
\theoremstyle{definition} \newtheorem{ex}[thm]{Example}
\theoremstyle{definition} \newtheorem{ej}[thm]{Exercise}

\newcommand*{\todo}[1]{\textcolor{red}{(TO DO: #1)}}

% Math Bold
\newcommand{\C}{\mathbb C}
\newcommand{\RP}{\mathbb{RP}}
\newcommand{\CP}{\mathbb{CP}}
\newcommand{\Q}{\mathbb Q}
\newcommand{\Z}{\mathbb Z}
\newcommand{\N}{\mathbb N}
\newcommand{\RR}{\mathbb R}
\newcommand{\bbA}{\mathbb A}
\newcommand{\bbT}{\mathbb T}
\newcommand{\bbR}{\mathbb R}
\newcommand{\bbK}{\mathbb K}
\newcommand{\bbS}{\mathbb S}
\newcommand{\bbH}{\mathbb H}
\newcommand{\bbB}{\mathbb B}
\newcommand{\bbX}{\mathbb X}
\newcommand{\bbE}{\mathbb E}
\newcommand{\bbP}{\mathbb P}
\newcommand{\bbI}{\mathbb I}

% Calligraphic
\newcommand{\cA}{\mathcal A}
\newcommand{\cB}{\mathcal B}
\newcommand{\cC}{\mathcal C}
\newcommand{\cD}{\mathcal D}
\newcommand{\cE}{\mathcal E}
\newcommand{\cF}{\mathcal F}
\newcommand{\cH}{\mathcal H}
\newcommand{\ch}{\mathcal h}
\newcommand{\cG}{\mathcal G}
\newcommand{\cI}{\mathcal I}
\newcommand{\cJ}{\mathcal J}
\newcommand{\cK}{\mathcal K}
\newcommand{\cL}{\mathcal L}
\newcommand{\cM}{\mathcal M}
\newcommand{\cN}{\mathcal N}
\newcommand{\cO}{\mathcal O}
\newcommand{\cP}{\mathcal P}
\newcommand{\cR}{\mathcal R}
\newcommand{\cS}{\mathcal S}
\newcommand{\cT}{\mathcal T}
\newcommand{\cU}{\mathcal U}
\newcommand{\cV}{\mathcal V}
\newcommand{\cW}{\mathcal W}
\newcommand{\cX}{\mathcal X}
\newcommand{\cY}{\mathcal Y}
\newcommand{\cZ}{\mathcal Z}

% Fraktur
\newcommand{\fA}{\mathfrak A}
\newcommand{\fB}{\mathfrak B}
\newcommand{\fC}{\mathfrak C}
\newcommand{\fD}{\mathfrak D}
\newcommand{\fH}{\mathfrak H}
\newcommand{\fP}{\mathfrak P}
\newcommand{\fS}{\mathfrak S}
\newcommand{\fT}{\mathfrak T}
\newcommand{\fU}{\mathfrak U}
\newcommand{\fV}{\mathfrak V}
\newcommand{\fX}{\mathfrak X}
\newcommand{\fR}{\mathfrak R}
\newcommand{\fY}{\mathfrak Y}
\newcommand{\fZ}{\mathfrak Z}
\newcommand{\fL}{\mathfrak L}
\newcommand{\ff}{\mathfrak f}
\newcommand{\fm}{\mathfrak m}
\newcommand{\fn}{\mathfrak n}
\newcommand{\fs}{\mathfrak s}
\newcommand{\ft}{\mathfrak t}
\newcommand{\fg}{\mathfrak{g}}
\newcommand{\fz}{\mathfrak{z}}
\newcommand{\fE}{\mathfrak E}
\newcommand{\fh}{\mathfrak h}
\newcommand{\fa}{\mathfrak a}
\newcommand{\fc}{\mathfrak c}

% Script
\newcommand{\sC}{\mathscr C}
\newcommand{\sX}{\mathscr X}
\newcommand{\sD}{\mathscr D}
\newcommand{\sU}{\mathscr U}
\newcommand{\sO}{\mathscr O}
\newcommand{\sH}{\mathscr H}
\newcommand{\sF}{\mathscr F}
\newcommand{\sL}{\mathscr L}
\newcommand{\sG}{\mathscr G}
\newcommand{\sJ}{\mathscr J}
\newcommand{\sS}{\mathscr S}

% Sans Serif
\newcommand{\sfA}{\mathsf A}
\newcommand{\sfB}{\mathsf B}
\newcommand{\sfC}{\mathsf C}
\newcommand{\sfD}{\mathsf D}
\newcommand{\sfE}{\mathsf E}
\newcommand{\sfF}{\mathsf F}
\newcommand{\sfG}{\mathsf G}

% Miscellaneous
\DeclareMathOperator{\pr}{pr}
\DeclareMathOperator{\GL}{GL}

\title{Differential Geometry}
\author{Guillermo Gallego}

\begin{document}
\begin{titlepage}
\setcounter{page}{1}
\thispagestyle{empty}
\begin{center}
\Huge \bfseries Differential Geometry \par\vspace{2em}
\Large Guillermo Gallego \par\bigskip
\end{center}

\vfill
\hrule \vspace{1em}
\begin{center}
\textsc{Departamento de Matemáticas, Universidad Autónoma de Madrid} \\ \vspace{1em}
\textit{Email:} \normalfont\href{mailto:guillermo.gallego@uam.es}{guillermo.gallego@uam.es} \hspace{1em}
\textit{URL:} \normalfont\href{https://guillegallego.xyz}{https://guillegallego.xyz}
\end{center}
\end{titlepage}

\tableofcontents
\chapter{Smooth manifolds}
\section{Definition and examples}
\subsection{Cartesian spaces}
\begin{defn}
By a \emph{Cartesian space} or \emph{abstract coordinate space} we mean an open subset $U\subset \mathbb{R}^n$, for some positive integer $n$, called the \emph{dimension of $U$}, $n=\dim U$. A map $f:U\subset \mathbb{R}^n \rightarrow V \subset \mathbb{R}^m$ between two Cartesian spaces is \emph{smooth} if all the iterated partial derivatives of $f$ exist and are continuous in $U$.
\end{defn}

The chain rule implies that, given Cartesian spaces $U\subset \mathbb{R}^n$, $V\subset \mathbb{R}^m$ and $W\subset \mathbb{R}^p$ and smooth maps $f:U\rightarrow V$ and $g:V\rightarrow W$, the composition $g\circ f:U\rightarrow W$ is also smooth. Therefore, there is a category, that we denote by $\mathbf{CartSp}$ whose objects are Cartesian spaces and whose morphisms are the smooth maps between them. An isomorphism in the category $\mathbf{CartSp}$ is a \emph{diffeomorphism}, by definition, a smooth bijective map whose inverse is also smooth.

\begin{defn}
Given a smooth map $f:U\subset \mathbb{R}^n\rightarrow V\subset \mathbb{R}^m$ between Cartesian spaces and a point $x\in U$, the \emph{differential of $f$ at $x$} is the linear map $d_xf:\mathbb{R}^n \rightarrow \mathbb{R}^m$ defined by the Jacobian matrix of $f$ at $x$
\begin{equation*}
  J_x f =
  \begin{pmatrix}
\frac{\partial f_1}{\partial x_1} (x) & \dots & \frac{\partial f_1}{\partial x_n} (x) \\
& \ddots & \\
\frac{\partial f_m}{\partial x_1} (x) & \dots & \frac{\partial f_m}{\partial x_n} (x)
  \end{pmatrix}.
\end{equation*}
Here, we are denoting by $f_i:U\rightarrow \mathbb{R}$ the composition of $f$ with the projection $\mathrm{pr}_i:V\rightarrow \mathbb{R}$ to the $i$-th coordinate.
\end{defn}

Again, the chain rule implies that the differential is ``functorial'' in the sense that, if $f:U\rightarrow V$ and $g:V\rightarrow W$ are smooth maps between Cartesian spaces, for any $x\in U$, we have
\begin{equation*}
d_x(g\circ f) = d_{f(x)} g \circ d_x f.
\end{equation*}
We can reformulate this as follows. Consider the category $\mathbf{CartSp}_*$ of \emph{based} Cartesian spaces, whose objects are pairs $(U\subset \mathbb{R}^n,x)$, formed by a Cartesian space $U$ and a point $x\in U$, and such that the morphisms $f:(U,x)\rightarrow (V,y)$ are given by smooth maps $f:U\rightarrow V$ such that $f(x)=y$. If we denote by $\mathbf{Vect}_{\mathbb{R}}$ the category of real vector spaces with linear maps, the chain rule is the statement that there is a functor, called the \emph{tangent functor}
\begin{equation*}
T: \mathbf{CartSp}_* \longrightarrow \mathbf{Vect}_{\mathbb{R}}
\end{equation*}
sending a based Cartesian space $(U,x)$ to the vector space $T_xU:= \mathbb{R}^{\dim U}$ and a map $f:(U,x)\rightarrow (V,y)$ to the linear map $d_xf:T_xU=\mathbb{R}^{\dim U} \rightarrow T_yV=\mathbb{R}^{\dim V}$. An immediate consequence is the invariance of dimension.

\begin{prop}[Invariance of dimension]
If two Cartesian spaces $U$ and $V$ are diffeomorphic, then $\dim U = \dim V$.
\end{prop}

\begin{proof}
Functors preserve isomorphisms. Therefore, if $f:U\rightarrow V$ is a diffeomorphism, then, for any $x\in U$, the linear map $d_xf:T_xU\rightarrow T_{f(x)}V$ is an isomorphism and thus $\dim U = \dim V$.
\end{proof}

The converse of the above is true only locally. This is the statement of the well-known inverse function theorem.

\begin{thm}[Inverse function theorem]
Let $f:U\rightarrow V$ be a smooth map of Cartesian spaces and let $x\in U$ be a point such that $d_xf:T_xU \rightarrow T_{f(x)}V$ is a linear isomorphism. Then, there exists some open neighbourhood $U'\subset U$ of $x$ such that the restriction $f|_{U'}:U'\rightarrow f(U')\subset V$ is a diffeomorphism.
\end{thm}

\subsection{Definition of smooth manifold}
\begin{defn}
Let $M$ be a topological space. A (local) \emph{Cartesian coordinate system} or \emph{chart} on $M$ is a tuple $(U,V,\varphi)$ formed by an open subset $U\subset M$, a Cartesian space $V\subset \mathbb{R}^n$, and a homeomorphism $\varphi:U\rightarrow V$. An \emph{atlas} on $M$ is a set of charts $\mathcal{A}$ such that $M=\bigcup_{(U,V,\varphi)\in \mathcal{A}} U$. We say that $M$ is \emph{locally Cartesian} if there exists an atlas on $M$.
\end{defn}

\begin{defn}
A topological space $M$ is a \emph{topological manifold} if it is
\begin{enumerate}
  \item locally Cartesian,
  \item Hausdorff, and
        \item second countable.
\end{enumerate}
\end{defn}

\begin{ej}
Can you think of examples of topological spaces which are locally Cartesian but not Hausdorff, or not second countable?
\end{ej}

Let $M$ be a topological space. Suppose that there are two charts $(U_1,V_1,\varphi_1)$ and $(U_2,V_2,\varphi_2)$ with $U_1\cap U_2 \neq \varnothing$. The homeomorphisms $\varphi_1$ and $\varphi_2$ restrict to homeomorphisms $\varphi_{12}:=\varphi_1|_{U_1\cap U_2}:U_1\cap U_2\rightarrow \varphi_1(U_1\cap U_2) \subset V_1$ and $\varphi_{21}:=\varphi_2|_{U_1\cap U_2}:U_1\cap U_2\rightarrow \varphi_2(U_1\cap U_2) \subset V_2$. We thus obtain a homeomorphism
\begin{equation*}
\psi_{12}:= \varphi_{21} \circ \varphi_{12}^{-1} : \varphi_1(U_1\cap U_2) \longrightarrow \varphi_2(U_1\cap U_2),
\end{equation*}
called the \emph{coordinate change map}.

\begin{defn}
Let $M$ be a topological manifold and let $\mathcal{A}$ be an atlas on $M$. We say that $\mathcal{A}$ is a \emph{smooth atlas} if, for every two charts $(U_1,V_1,\varphi_1), (U_2,V_2,\varphi_2)\in \mathcal{A}$, the coordinate change map $\psi_{12}: \varphi_1(U_1\cap U_2) \rightarrow \varphi_2(U_1\cap U_2)$ is a diffeomorphism.
\end{defn}

One could initially try to define a smooth manifold as a topological manifold endowed with a smooth atlas $\mathcal{A}$. However, two different atlases can determine essentially the same ``smooth structure''. Two smooth atlases $\mathcal{A}_1$ and $\mathcal{A}_2$ on a topological manifold $M$ are \emph{compatible} if their union $\mathcal{A}_1 \cup \mathcal{A}_2$ is also a smooth atlas on $M$. In other words, if for any two charts $(U_1,V_1,\varphi_1) \in \mathcal{A}_1$ and $(U_2,V_2,\varphi_2) \in \mathcal{A}_2$, the coordinate change map $\psi_{12}$ is a diffeomorphism. Given any smooth atlas $\mathcal{A}$ on $M$, we can consider the union of all the smooth atlas compatible with it. This union is again an atlas $\bar{\mathcal{A}}$, which is maximal with respect to inclusion. Such a maximal atlas is called a \emph{smooth structure} on $M$. If two smooth manifolds have the same maximal atlas, then they are essentially the same smooth manifold.

\begin{defn}
A \emph{smooth manifold} is a pair $(M,\mathcal{A})$ formed by a topological manifold $M$ and a smooth structure $\bar{\mathcal{A}}$ on $M$.
\end{defn}

If $(M,\mathcal{A})$ is a smooth manifold and $(U_1,V_1,\varphi_1)\in \mathcal{A}$ is a chart, for every point $x\in U$, we define the \emph{dimension of $M$ at $x$} to be the dimension
\begin{equation*}
\dim_x M:= \dim V_1.
\end{equation*}
Note that this number does not depend on the choice of chart, since, for any other chart $(U_2,V_2,\varphi_2)$, the coordinate change map $\psi_{12}$ induces a linear isomorphism $$d_{\varphi_1(x)}\psi_{12}:T_{\varphi_1(x)}V_1 \rightarrow T_{\varphi_2(x)}V_2.$$ Therefore, we obtain a well defined map
\begin{equation*}
\dim: M \rightarrow \mathbb{N}, x \mapsto \dim_x M.
\end{equation*}
Moreover, note that, since, for every chart $(U,V,\varphi)\in \mathcal{A}$ and for every point $x\in U$, the open subset $U\subset M$ is contained in the preimage $\dim^{-1}(\dim_x M)$, the map $\dim$ is continuous if we endow $\mathbb{N}$ with the discrete topology. Thus the dimension is constant on connected components and it makes sense to talk about the \emph{dimension} $\dim M$ of a connected smooth manifold $M$.

A connected smooth manifold of dimension $n$ is called a (smooth) \emph{$n$-manifold} or \emph{$n$-fold}. A $1$-fold is also called a (smooth) \emph{curve}. A $2$-fold is called a (smooth) \emph{surface}.

\subsection{Smooth maps}
Let $f:M\rightarrow N$ be a continuous map between two topological spaces. Suppose that there are charts $(U,V,\varphi)$ and $(U',V',\varphi')$ with $f(U)\subset U'$. The \emph{localization} of $f$ with respect to the charts $(U,V,\varphi)$ and $(U',V',\varphi')$ is the map $\tilde{f}: V \rightarrow V'$ given by the diagram
\begin{center}
  \begin{tikzcd}
   U \ar{r}{f|_U} \ar{d}{\varphi} & U' \ar{d}{\varphi'} \\
   V \ar{r}{\tilde{f}} & V'.
  \end{tikzcd}
\end{center}
That is, $\tilde{f}=\varphi'\circ f \circ \varphi^{-1}$.

\begin{defn}
A map $f:M\rightarrow N$ between two smooth manifolds is \emph{smooth} if, for every two charts $(U,V,\varphi)$ and $(U',V',\varphi')$ of $M$ and $N$ with $f(U)\subset U'$, the localization $\tilde{f}:V\rightarrow V'$ is a smooth map.
\end{defn}

Note that localization commutes with composition. That is, if $f:M\rightarrow N$ and $g:N\rightarrow P$ are maps between smooth manifolds and $(U,V,\varphi)$, $(U',V',\varphi')$ and $(U'',V'',\varphi'')$ are charts of $M$, $N$ and $P$, respectively, with $f(U)\subset U'$ and $g(U')\subset U''$, then the localization $\widetilde{(g\circ f)}$ of the composition of $f$ and $g$ coincides with the composition of the localizations $\tilde{g}\circ \tilde{f}$. Therefore, it follows from the chain rule that if $f$ and $g$ are smooth maps then $g\circ f$ is also a smooth map.

It follows that smooth manifolds with smooth maps form a category, that we denote by $\mathbf{SmoothMan}$. Isomorphisms in this category are also called \emph{diffeomorphisms}. By definition, a diffeomorphism is a smooth bijective map whose inverse is also smooth.

\subsection{Some examples}
\subsubsection{Cartesian spaces} A Cartesian space $U$ is trivially a smooth manifold if we simply endow it with the atlas $\mathcal{A}=\left\{ (U,U,\mathrm{id}_{U}) \right\}$. By construction, maps between Cartesian spaces are smooth if and only if they are smooth as maps between smooth manifolds. Therefore we have a functor $\mathbf{CartSp}\rightarrow \mathbf{SmoothMan}$, which is clearly fully faithful.

\subsubsection{Spheres} The \emph{$n$-dimensional unit sphere} is the space
\begin{equation*}
\mathbb{S}^n=\left\{(x_1,\dots,x_n,x_{n+1}) \in \mathbb{R}^{n+1}: x_1^2+\dots +x_n^2+x_{n+1}^2=1  \right\}.
\end{equation*}
Let $p_N=(0,\dots,0,1)$ and $p_S=(0,\dots,0,-1)$ denote the north and south poles of the sphere, and consider the open subsets $U_1=\mathbb{S}^n \setminus \left\{ p_N \right\}$ and $U_2=\mathbb{S}^n \setminus \left\{ p_S \right\}$. The stereographic projections $\varphi_1:U_1\rightarrow \mathbb{R}^n$ and $\varphi_1:U_2\rightarrow \mathbb{R}^n$, centered on the south and north poles, respectively, are homeomorphisms, and one checks easily that the coordinate change is a diffeomorphism. Therefore, the set $\mathcal{A}=\left\{ (U_1,\mathbb{R}^n,\varphi_1), (U_2,\mathbb{R}^n,\varphi_2) \right\}$ is an smooth atlas, which endows $\mathbb{S}^n$ with the structure of an $n$-dimensional smooth manifold.

\subsubsection{Projective spaces} Consider the real projective space $\mathbb{RP}^n$. It is covered by subsets of the form
\begin{equation*}
D(x_i)= \left\{ (x_0:x_1:\dots:x_n) : x_i \neq 0 \right\},
\end{equation*}
which are open subsets homeomorphic to $\mathbb{R}^n$ through the map
\begin{equation*}
\varphi_i: D(x_i)\rightarrow \mathbb{R}^n: (x_0:\dots:x_i:\dots:x_n) \mapsto \left(\frac{x_0}{x_i},\frac{x_1}{x_i},\dots,\frac{x_n}{x_i}\right).
\end{equation*}
The coordinate changes $\psi_{ij}:D(x_i)\cap D(x_j) \rightarrow D(x_i)\cap D(x_j)$ are of the form
\begin{equation*}
\psi_{ij}\left(\frac{x_0}{x_i},\frac{x_1}{x_i},\dots,\frac{x_n}{x_i}\right) = \left(\frac{x_0}{x_j},\frac{x_1}{x_j},\dots,\frac{x_n}{x_j}\right),
\end{equation*}
so they are clearly diffeomorphisms.

Any linear isomorphism $\hat{f}:\mathbb{R}^{n+1}\rightarrow \mathbb{R}^{n+1}$ induces a map of projective spaces $f:\mathbb{RP}^n\rightarrow \mathbb{RP}^n$ (a \emph{homography}). As a map between manifolds, $f$ is smooth. Indeed, if we restrict $f$ to the complement $D(H)$ of any hyperplane $H\subset \mathbb{RP}^n$, the induced localization $\tilde{f}:\mathbb{R}^n \cong D(H)\rightarrow D(f(H))\cong \mathbb{R}^n$ is an affine map, and thus smooth.

\subsubsection{Different smooth structures on the same topological manifold} One can endow a topological manifold with two smooth structures which are essentially different; that is, with two different maximal atlases. In order to do this, we just need to find two atlases which are not compatible. For example, we can consider the homeomorphisms
\begin{equation*}
\mathrm{id}_{\mathbb{R}}: \mathbb{R} \rightarrow \mathbb{R}, t \mapsto t,
\end{equation*}
and
\begin{equation*}
\varphi: \mathbb{R} \rightarrow \mathbb{R}, t \mapsto t^3.
\end{equation*}
These determine two different atlases $\mathcal{A}_1=\left\{ (\mathbb{R},\mathbb{R},\mathrm{id}_{\mathbb{R}}) \right\}$ and $\mathcal{A}_2=\left\{ (\mathbb{R},\mathbb{R},\varphi) \right\}$ on $\mathbb{R}$ which are not compatible, since the coordinate change map $\psi(t)=t^3$ is not a diffeomorphism (the inverse $\psi^{-1}(t)=t^{1/3}$ is not smooth). Therefore, the associated maximal atlases $\bar{\mathcal{A}}_1$ and $\bar{\mathcal{A}}_2$ are different, and thus they determine two different smooth structures.

However, in this case the smooth manifolds $(\mathbb{R},\bar{\mathcal{A}}_1)$ and $(\mathbb{R},\bar{\mathcal{A}}_2)$ are diffeomorphic. Indeed, the map $\varphi(t)=t^3$ determines a diffeomorphism between $(\mathbb{R},\bar{\mathcal{A}}_1)$ and $(\mathbb{R},\bar{\mathcal{A}}_2)$, since its localization $\tilde{f}$ with respect to the charts $(\mathbb{R},\mathbb{R},\mathrm{id}_{\mathbb{R}})$ and $(\mathbb{R},\mathbb{R},\varphi)$ is just the identity.

It turns out that, up to diffeomorphism, $\mathbb{R}$ only admits one smooth structure, but it is important to know that this is not the general situation. Indeed, there are some topological manifolds which admit several different smooth structures which are not diffeomorphic to each other.

\subsubsection{The classification problem of smooth manifolds}
A natural question is then to classify smooth manifolds up to diffeomorphism. The problem is solved in low dimension. Let $M$ be a connected smooth manifold.
\begin{itemize}
\item If $\dim M = 1$, then $M$ is either diffeomorphic to a circle $\mathbb{S}^1$ (if it is compact) or to the real line $\mathbb{R}$ (if it is not compact). (See \todo{CITE}).
\item In dimension $2$, every topological manifold admits a unique smooth structure, and thus the classification is reduced to the classification of topological surfaces (see \todo{CITE} for that classification). We conclude that, if $\dim M =2$, then it is diffeomorphic to a sphere with $g$ handles $\Sigma_g$ (if it is orientable) or to a sphere with $k$ crosscaps $X_k$ (if it is not orientable).
\item In dimension $3$ the classification also coincides with the topological classification (see \todo{CITE}), but the topological classification is not complete.
\end{itemize}

In dimension $\geq 4$ (and specifically in dimension $4$) is where the problem reaches its full complexity, since one can find topological manifolds which admit several smooth structures which are not diffeomorphic to each other. If $M$ is compact and $\dim M \geq 5$, then $M$ can admit several smooth structures not diffeomorphic to each other, but only a finite amount. For example, the sphere $\mathbb{S}^7$ admits $28$ smooth structures not diffeomorphic to each other (the $27$ smooth structures different from the ``natural'' one are called the \emph{exotic spheres}). A $4$-dimensional compact smooth manifold can admit an infinite (countable) amount of non-diffeomorphic smooth structures. In the $4$-dimensional non-compact case there are really surprising results. For example, while $\mathbb{R}^n$ admits a unique smooth structure up to diffeomorphism for $n\neq 4$, the space $\mathbb{R}^4$ admits an infinite (\emph{uncountable}) amount of non-diffeomorphic smooth structures.

\subsection{Partitions of unity and the coordinate ring}
\begin{defn}
  A \emph{smooth partition of unity} on a smooth manifold $M$ is a family $\Theta$ of smooth functions $\theta:M\rightarrow \mathbb{R}$ such that
  \begin{enumerate}
    \item $\theta \geq 0$ for every $\theta \in \Theta$,
    \item it is \emph{locally finite}, meaning that for each point $x\in M$ there exists an open neighbourhood $U\subset M$ of $x$ such that all but a finite amount of the $\theta \in \Theta$ vanish at $U$, and
          \item $\sum_{\theta \in \Theta} \theta = 1$.
  \end{enumerate}

  An \emph{open cover} $\mathfrak{U}$ of $M$ is a family of open subsets $U\subset M$ such that $M=\bigcup_{U\in \mathfrak{U}} U$. A smooth partition of unity \emph{subordinate to $\mathfrak{U}$} is a smooth partition of unity $\Theta$ together with a bijection $\mathfrak{U}\rightarrow \Theta, U\mapsto \theta_U$ such that, for each $U\in \mathfrak{U}$, the function $\theta_U$ vanishes away from $U$.
\end{defn}

\begin{defn}
  Let $M$ be a smooth manifold.
  \begin{itemize}
    \item Let $A\subset M$ be a closed subset and let $U\subset M$ be an open neighbourhood of $A$. A \emph{bump function} on $A$ is a smooth non-negative function $\theta:M\rightarrow \mathbb{R}$ with $\theta|_A = 1$ and $\theta|_{M\setminus U} = 0$.
          \item Let $A$ and $B$ be two disjoint closed subsets of $M$. An \emph{Uryshon separating function} of $A$ and $B$ is a smooth non-negative function $\theta:M\rightarrow \mathbb{R}$ with $\theta|_A=1$ and $\theta|_B = 0$.
  \end{itemize}
\end{defn}

The existence of smooth bump functions in $\mathbb{R}^n$ implies the existence of bump functions, Uryshon separating functions and smooth partitions of unity for any smooth manifold. This is one of the most important features of differential geometry. Notably, there are no counterparts to these notions in analytic or algebraic geometry. Partitions of unity are a crucial tool in differential geometry, since they allow one to construct global objects from objects which are only locally defined. The clearest example is the following.

\begin{prop}[Tietze's extension theorem]
Let $M$ be a smooth manifold, $A\subset M$ a closed subset, $U$ an open neighbourhood of $A$ and $f:U\rightarrow \mathbb{R}$ a smooth function. Then there exists a smooth function $F:M\rightarrow \mathbb{R}$ such that $F|_A = f$.
\end{prop}

\begin{proof}
Consider the open cover $\mathfrak{U}=\left\{ U, M\setminus A \right\}$ and let $\Theta=\left\{ \theta_U,\theta_{M\setminus A} \right\}$ be a smooth partition of unity subordinate to $\mathfrak{U}$. We can then define $F=f \theta_U + \theta_{M \setminus A}$, which clearly coincides with $f$ at any point of $A$.
\end{proof}

Another application of the existence of smooth partitions of unity is the existence of proper functions.

\begin{prop}
For any smooth manifold $M$ there exists a proper map $f:M\rightarrow \RR$.
\end{prop}

\begin{proof}
\todo{}
\end{proof}

The \emph{coordinate ring} of a smooth manifold $M$ is the $\RR$-algebra $C^\infty(M)$ of functions $M\rightarrow \RR$. With any point $x\in M$ we can associate the form
\begin{equation*}
\bar{x}:C^\infty(M)\rightarrow \RR, f \mapsto f(x).
\end{equation*}
It turns out that the existence of smooth partitions of unity allows us to recover all the information about the smooth structure of $M$ from the coordinate ring $C^\infty(M)$. Indeed, we have the following.

\begin{thm}
The map $M \rightarrow \mathrm{Hom}_{\mathbb{R}}(C^\infty(M),\mathbb{R}), x\mapsto \bar{x}$ is a bijection.
\end{thm}
\begin{proof}
\todo{}
\end{proof}

\begin{corol}
The functor $C^\infty:\mathbf{SmoothMan}\rightarrow \mathbf{Alg}_{\mathbb{R}}^{\mathrm{op}}$ sending a smooth manifold $M$ to its coordinate ring $C^\infty(M)$ and a smooth map $f:M\rightarrow N$ to the morphism of $\mathbb{R}$-algebras
\begin{equation*}
f^*: C^\infty(N) \rightarrow C^\infty(M), \varphi \mapsto \varphi \circ f,
\end{equation*}
is fully faithful.
\end{corol}
\begin{proof}
\todo{}
\end{proof}

\subsection{Quotients}
??

\subsection{Lie groups}
??

\subsection{Vector bundles}
\begin{defn}
  Let $M$ be a topological space and let $\mathbb{K}$ be either $\mathbb{R}$ or $\mathbb{C}$. A \emph{$\mathbb{K}$-vector bundle of rank $r$ over $M$} is a topological space $E$ together with a surjective continuous map $\pi:E\rightarrow M$ such that:
  \begin{itemize}
    \item For each $x\in M$, the fibre $E_x:=\pi^{-1}(x)$ is endowed with the structure of an $r$-dimensional $\mathbb{K}$-vector space.
    \item For each $x\in M$ there exists a neighbourhood $U$ of $x$ in $M$ and a homeomorphism $\varphi: E|_U:= \pi^{-1}(U)\rightarrow U \times \mathbb{K}^r$ (called a \emph{local trivialization of $E$ over $U$}) such that
          \begin{itemize}
            \item If $\mathrm{pr}_U: U \times \mathbb{K}^r \rightarrow U$ denotes the natural projection, the following diagram commutes
                  \begin{center}
                    \begin{tikzcd}
                      E|_{U} \arrow{rd}{\pi} \rar{\varphi} & U \times \mathbb{K}^r \dar{\mathrm{pr}_U} \\
                      & U.
                    \end{tikzcd}
                  \end{center}
            \item For each $x\in U$, the restriction of $\varphi|_{E_x}: E_x \rightarrow \left\{ x \right\}\times \mathbb{K}^r \cong \mathbb{K}^r$ is a $\mathbb{K}$-vector space isomorphism.
          \end{itemize}
  \end{itemize}
  When $\mathbb{K}=\mathbb{R}$, we say that $E$ is a \emph{real} vector bundle, while if $\mathbb{K}=\mathbb{C}$ we say that $E$ is a \emph{complex} vector bundle. In this text, when not specified, by a vector bundle we just mean a real vector bundle. A vector bundle of rank $1$ (real or complex) is called a \emph{line bundle}.

  The space $E$ is called the \emph{total space} of the bundle, $M$ is called its \emph{base}, and $\pi$ is the \emph{projection}. We typically introduce a vector bundle by saying ``$E$ is a vector bundle over $M$'', or ``$E\rightarrow M$ is a vector bundle''.

  If $E$ and $M$ are smooth manifolds, $\pi$ is a smooth map, and the local trivializations can be chosen to be diffeomorphisms, then $E$ is called a \emph{smooth vector bundle}.

  A smooth local trivialization (resp. a continuous local trivialization) defined over all of $M$ is called a \emph{global trivialization} (resp. a \emph{continuous global trivialization}). If such a global trivialization exists, then $E$ itself is diffeomorphic (resp. homeomorphic) to $M\times \mathbb{K}^r$ and $E$ is said to be a \emph{trivial bundle} (resp. a \emph{topologically trivial bundle}).
\end{defn}

\subsubsection{The Möbius bundle}
The simplest example of a non-trivial vector bundle is given by the well-known Möbius construction. We start by defining the following equivalence relation on $\RR^2$
\begin{equation*}
(x,y) \sim (x' ,y' ) \text{ if and only if } (x' ,y' )= (x+n, (-1)^n y) \text{ for some } n \in \Z.
\end{equation*}
We let $E=\RR^2/\sim$ be the quotient space and denote by $q:\RR^2 \rightarrow E$ the quotient map. To help visualize the construction, note that the image under $q$ of a rectangle $[0,1]\times [-l,l]$ is the famous Möbius band.

Consider now the projection to the first factor $\mathrm{pr}_1:\RR^2\rightarrow \RR$ and the covering map $\epsilon:\RR\rightarrow \mathbb{S}^1$, $\epsilon(t)=\exp(2\pi i t)$. Since $\epsilon \circ \mathrm{pr}_1$ is constant on each equivalence class, it descends to a map $\pi:E\rightarrow \mathbb{S}^1$ making the following diagram commute
\begin{center}
  \begin{tikzcd}
   \RR^2 \ar{r}{q} \ar{d}{\mathrm{pr}_1} & E \ar{d}{\pi} \\
   \RR \ar{r}{\epsilon} & \mathbb{S}^1.
  \end{tikzcd}
\end{center}
The space $E$ admits a unique smooth structure such that $q$ is a smooth covering map and $\pi:E\rightarrow \mathbb{S}^1$ is a smooth real line bundle.

\begin{prop}
The Möbius bundle $E\rightarrow \mathbb{S}^1$ is not trivial (i.e. it does not admit a global trivialization).
\end{prop}

\begin{proof}
  Suppose that there is a global trivialization $\varphi:E\rightarrow \mathbb{S}^1 \times \mathbb{R}$ and consider the map
  \begin{equation*}
s:\mathbb{S}^1 \rightarrow \mathbb{S}^1  \times \mathbb{R}, x \mapsto (x,1),
  \end{equation*}
  and the composition $s' = \varphi^{-1}\circ s: \mathbb{S}^1\rightarrow E$.
  Since $\varphi$ is a linear isomorphism on each fibre, for each $x\in \mathbb{S}^1$, the point $s'(x)\in E_x$ is a non-zero element of the vector space $E_x$. The map $s'$ lifts to a map $f:\RR \rightarrow \RR^2$ in the sense that
  \begin{equation*}
s'(\exp(2\pi i t)) = q(f(t)).
  \end{equation*}
  This map must satisfy the conditions $f(t)\neq 0$ for every $t\in \mathbb{R}$ and $f(1)=-f(0)$, contradicting Bolzano's theorem.
\end{proof}

\subsubsection{Transition functions}
\begin{prop}
  Let $\pi:E\rightarrow M$ be a smooth $\mathbb{K}$-vector bundle of rank $r$ over $M$. Suppose that $\varphi:E|_U \rightarrow U \times \mathbb{K}^r$ and $\psi:E|_V \rightarrow V \times \mathbb{K}^r$ are two smooth local trivializations of $E$ with $U\cap V \neq \varnothing$. Then there exists a smooth map $\tau:U\cap V \rightarrow \GL_r(\bbK)$, called the \emph{transition function}, such that the composition $\varphi \circ \psi^{-1}:(U\cap V) \times \bbK^r \rightarrow (U\cap V) \times \bbK^r$ has the form
  \begin{equation*}
\varphi \circ \psi^{-1}(x,v) = (x, \tau(x)v).
  \end{equation*}
\end{prop}

\begin{proof}
We have the following commutative diagram
\begin{center}
  \begin{tikzcd}
(U\cap V) \times \bbK^r \ar{dr}{\pr_{U\cap V}} & \pi^{-1}(U\cap V) \ar{l}{\psi} \ar{r}{\varphi} \ar{d}{\pi} & (U\cap V) \times \bbK^r \ar{dl}{\pr_{U\cap V}} \\
& U \cap V. &
  \end{tikzcd}
\end{center}
It follows that
\begin{equation*}
\varphi \circ \psi^{-1}(x,v) = (x,\sigma(x,v))
\end{equation*}
for some smooth map $\sigma: (U\cap V) \times \bbK^r \rightarrow \bbK^r$. Moreover, for each $x\in U\cap V$, the map $v \mapsto \sigma(x,v)$ is a linear automorphism of $\bbK^r$, and thus we can write $\sigma(x,v)=\tau(x) v$, for some $\tau(x) \in \GL_r(\bbK)$. Since $\sigma$ is smooth, the map $x \mapsto \tau(x)$ must also be smooth.
\end{proof}

It turns out that a vector bundle can be reconstructed from its transition functions.

\begin{prop}
  Let $M$ be a smooth manifold and suppose that, for each $x\in M$, we are given an $r$-dimensional $\bbK$-vector space $E_x$, for some fixed positive integer $r$. Consider the disjoint union $E=\coprod_{x\in M} E_x$ and let $\pi:E\rightarrow M$ be the map sending each element of $E_x$ to the point $x$. Suppose furthermore that we are given the following data:
  \begin{enumerate}
    \item an open cover $\mathfrak{U}$ of $M$,
          \item for each $U\in \mathfrak{U}$, a bijection $\varphi_U:\pi^{-1}(U)\rightarrow U \times \bbK^r$ whose restriction to each $E_x$ is a $\bbK$-vector space isomorphism between $E_x$ and $\left\{ x \right\}\times \bbK^r \cong \bbK^r$.,
          \item for each $U,V \in \mathfrak{U}$ with $U\cap V\neq \varnothing$, a smooth map $\tau_{UV}:U\cap V \rightarrow \GL_r(\bbK)$ such that the map $\varphi_U \circ \varphi_V^{-1}:(U\cap V)\times \bbK^r\rightarrow (U\cap V)\times \bbK^r$ has the form
          \begin{equation*}
\varphi_U \circ \varphi_V^{-1}(x,v) = (x,\tau_{UV}(x) v).
          \end{equation*}
  \end{enumerate}
  Then $E$ has a unique topology and smooth structure making it into a smooth manifold and a smooth rank $r$ $\bbK$-vector bundle over $M$, with $\pi$ as a projection and the $(U,\varphi_U)$, for $U\in \mathfrak{U}$, as smooth local trivializations.
\end{prop}

\begin{proof}
a
\end{proof}

\subsubsection{Constructions}

\subsubsection{Sections}
definition, extensions, module structure and reconstruction

\section{The tangent functor}
\subsection{Derivations and tangent space}
\subsection{Differential of a smooth map. The chain rule}
\subsection{The tangent bundle}
\subsection{Embeddings}
\subsection{Embedding theorem}

\section{Tensor calculus. Fields and forms}
\subsection{Multilinear algebra}

\chapter{Some geometric structures}
\section{Riemannian metrics}
\section{Symplectic structures}
\section{Almost complex structures}
\section{Hermitian metrics. Kähler structures}

\chapter{Differential equations on smooth manifolds}
\section{Integrating the flow of a vector field}
\section{Hamiltonian mechanics}
\section{Distributions and integrability}
\section{The Frobenius theorem}

\chapter{Lie groups}

\chapter{Cohomology}
\section{Chains and singular (co)homology}
\section{Integration on chains and the Stokes theorem}
\section{de Rham cohomology}
\section{The Poincaré lemma}
\section{The de Rham theorem}
\section{Mapping degree theory}

\chapter{Harmonic forms}
\section{The Hodge star operator and the Laplacian}
\section{The Hodge theorem}
\section{Applications of the Hodge theorem}

\chapter{Connections on vector bundles}
\section{The Levi-Civita connection}

\chapter{Connections on principal bundles}

%\nocite{*}
%\printbibliography
\end{document}
